Robots that cross pedestrian areas rely on trajectory prediction, yet
predictors are usually judged by displacement error rather than by how
safely a robot drives with them. We build a complete navigation stack in
MATLAB with a classical and a learned option at each layer: seven
predictors (constant velocity, Kalman, Social Force, MLP, LSTM, temporal
CNN and a social-grid CNN), split and adaptive conformal uncertainty, 2-D
and space-time A*, and pure pursuit, dynamic-window and sampling-based MPC
controllers. We evaluate 23 configurations on the same 120 episodes from
the ETH/UCY scenes, using a leak-free split and a reactive Social Force
crowd that follows the recorded routes. Prediction halves the contacts of
a space-time planner (1.47 to 0.67 per episode), and space-time A* with
MPC has 54\% fewer contacts than A* with pure pursuit. However, the learned
predictors are up to 10\% more accurate but no safer than constant
velocity: no paired difference is significant. Conformal margins lose
their 90\% coverage under the shift to an unseen scene and, used as hard
constraints, cut goal-reaching from 93\% to 55\% without reducing
contacts. The planner, not the predictor, is the lever for safety.
